\documentclass[../../../include/open-logic-section]{subfiles}

\begin{document}

\olfileid{sth}{z}{nat}
\olsection{आपल्या नैसर्गिक संख्यांची निवड}

%Let us take it for granted that the informal principle \stagesinf{}
%can be justified. It is also, though, worth commenting on the
%particular axiom we extracted from that informal principle.

\olref[infinity-again]{defnomega} मध्ये आपण ``नैसर्गिक संख्या'' या शब्दप्रयोगाची स्पष्ट व्याख्या केली. हा निर्णय कसा समजायचा? तो आधिभौतिक दावा नाही; काही संचांना नैसर्गिक संख्या म्हणून \emph{वागवण्याचा} तो निर्णय आहे. असे मानण्याची कारणे आपण \olref[sfr][rel][ref]{sec}, \olref[sfr][arith][ref]{sec} आणि \olref[sfr][infinite][dedekindsproof]{sec} मध्ये पाहिली. पण ती आणखी नेमकी सांगता येतील.

आपले अनंततेचे स्वयंसिद्धक \citet{VonNeumann1925} यांच्या पद्धतीनुसार आहे. मात्र त्याऐवजी पुढील स्वयंसिद्धकही स्वीकारता आले असते:

\begin{defish}
\emph{झर्मेलो यांचे \citeyear{Zermelo1908Untersuchungen} मधील अनंततेचे स्वयंसिद्धक.} असा एक संच $A$ आहे की $\emptyset \in A$ आणि $(\forall x \in A)\{x\} \in A$.
\end{defish}

आपण फोन न्यूमन यांच्या धर्तीवरचे अनंततेचे स्वयंसिद्धक न घेता झर्मेलो यांचे घेतले असते, तरी डेडेकिंड-अनंत संच आणि म्हणून डेडेकिंड बीजसंरचना मिळाली असती. झर्मेलो यांच्या पद्धतीत आपल्या बीजसंरचनेचा निवडक घटक पुन्हा $\emptyset$ हाच असता ($0$ चे आपले प्रतिरूप); पण एकास-एक फलन $x \mapsto x \cup \{x\}$ ऐवजी $x \mapsto \{x\}$ असे असते. याचा सर्वांत सरळ परिणाम असा की झर्मेलो यांच्या पद्धतीत $2$ म्हणजे $\{\{\emptyset\}\}$, तर आपण फोन न्यूमन यांच्यासह $2$ म्हणजे $\{\emptyset, \{\emptyset\}\}$ मानतो.

अनंततेच्या या दोन स्वयंसिद्धकांपैकी एकच का निवडावे? मुख्य व्यावहारिक कारण असे की फोन न्यूमन यांची पद्धत अनंतपार संख्यांसाठी चांगली विस्तारते. \olref[ordinals][]{chap} पासून हे पाहू. मात्र केवळ \emph{अंकगणित करण्याच्या} दृष्टीने दोन्ही पद्धती तितक्याच उपयोगी आहेत. म्हणून कोणी नैसर्गिक संख्या \emph{संचच आहेत} असे सांगितले, तर लगेच प्रश्न पडतो: \emph{त्या नेमक्या कोणत्या संच आहेत?}

हा नेमका प्रश्न \citet{Benacerraf1965} यांनी प्रसिद्ध केला. पण आपण यापूर्वी अनेकदा पाहिलेल्या व्यापक बाबीचे ते केवळ सर्वांत प्रसिद्ध उदाहरण आहे. मूलभूत मुद्दा असा: संच उपपत्तीच्या साहाय्याने अनेक ``सहजगत्या'' मानल्या जाणाऱ्या प्रकारच्या वस्तूंचे \emph{अनुकरण} करता येते: वास्तव संख्या, परिमेय संख्या, पूर्णांक आणि नैसर्गिक संख्या; तसेच क्रमित जोड्या, फलने आणि संबंध. पण संच उपपत्ती अशा अनुकरणाची \emph{एकमेव} निवड देत नाही. त्याच कामाला सरळपणे तितकेच उपयोगी पडले असते असे पर्याय \emph{नेहमीच} असतात.

\end{document}
